\documentclass[11pt]{article}
\usepackage[margin=1in]{geometry}
\usepackage{amsmath,amssymb,amsthm}
\usepackage[colorlinks=true,linkcolor=blue,citecolor=blue,urlcolor=blue]{hyperref}
\newtheorem{theorem}{Theorem}
\newtheorem{lemma}{Lemma}
\newcommand{\Fp}{\mathbb F_p}
\newcommand{\one}{\mathbf1}
\newcommand{\E}{\mathbb E}
\newcommand{\Prob}{\mathbb P}
\newcommand{\rank}{\operatorname{rank}}
\newcommand{\nullity}{\operatorname{nullity}}
\title{Odd block sizes: an addendum to the RSBM kernel argument}
\author{Research session note}
\date{November 15, 2026}
\begin{document}
\maketitle
\begin{abstract}
A standard vertex-insertion construction and a conditional finite-field calculation extend the companion even-block theorem to odd block sizes for every fixed even internal degree at least six. The cross graph remains arbitrary, prescribed independently of the internal graphs. Neither an effective probability rate nor an odd-block internal-degree-four result is obtained.
\end{abstract}

\begin{theorem}\label{thm:odd}
Fix even $k\ge6$ and $0<\ell<k$, and put $\delta=k-\ell$. Let $N$ tend to infinity through odd integers. For independent uniform simple $k$-regular internal adjacency matrices $C,D$ on $N$ vertices, and any prescribed simple $\ell$-biregular cross matrix $B$, put
\[
 A_B=\begin{pmatrix}C&B\\B^T&D\end{pmatrix}.
\]
Then
\[
 \sup_B\Prob\bigl(\nullity_{\mathbb R}(A_B-\delta I)\ge2\bigr)\longrightarrow0.
\]
\end{theorem}
The companion even-block theorem and admissibility for odd internal degrees therefore cover every admissible size for $d_1=3$ and every $d_1\ge5$. The uncovered cases are $(2,1)$ and odd block sizes for $d_1=4$.

\section{A standard vertex insertion and its distribution}
Write $n=N-1$, which is even, and $s=k/2$. On a simple $k$-regular graph $G$ on $[n]$, choose uniformly a set $E$ of $s$ vertex-disjoint edges. Delete these edges, insert the fixed root $N$, and join it to all $2s=k$ endpoints. This gives a simple $k$-regular graph $H$ on $[N]$. This odd-order construction is known; see also the odd-order argument in \cite{hlmpw}. We give the counting justification needed here.

Let $m=kn/2$. Uniformly over loopless bounded-degree multigraphs, counting edge copies separately, the number of $s$-edge matchings is
\begin{equation}\label{eq:matchcount}
 \frac{m^s}{s!}\bigl(1+O(n^{-1})\bigr).
\end{equation}
Indeed, after any fixed number of chosen disjoint edges only $O_k(1)$ edge copies meet their endpoints; thus the ordered count is $m^s(1+O(n^{-1}))$.

For a given simple $H$, a preimage is a pairing of its $k$ root-neighbors into pairs that are nonedges of $H-N$. Let $h(H)$ count these valid pairings, and set $L=(k-1)!!$. By \eqref{eq:matchcount}, the output law from uniform simple $G$ has mass
\[
 c_n h(H)(1+O(n^{-1}))
\]
with a uniform relative error where $h(H)>0$. Also $0\le h(H)\le L$, and $h(H)=L$ whenever the root belongs to no triangle. Under uniform simple $k$-regular $H$, the root-triangle probability is $O(n^{-1})$: expose its fixed-degree neighborhood in the configuration model and then condition on simplicity, an event of probability bounded below.

If $\mu_N$ denotes uniform simple $H$, the probability law with density $h/\E_{\mu_N}h$ differs from $\mu_N$ by $O(n^{-1})$ in total variation, since $\E_{\mu_N}h=L(1-O(n^{-1}))$. Normalizing the displayed output mass shows that it too differs from this density law by $O(n^{-1})$. Thus the surgery output is within $O(n^{-1})$ total variation of uniform simple $k$-regular $H$.

\section{Colored sampling and independent untouched matchings}
First generate $G$ as the sum of $k$ independent uniform perfect matchings $M_1,\ldots,M_k$ on $[n]$, allowing parallel edges with distinct colors. Sample $s$ independent colors $c_j$ uniformly from $[k]$ and vertices $v_j$ uniformly from $[n]$, choosing oriented edge $(v_j,M_{c_j}v_j)$. Let $V$ mean that all these edges are vertex-disjoint.

Uniformly in $G$, $\Prob(V\mid G)=1-O(n^{-1})$. Conditional on $G,V$, the unordered edges are exactly a uniform $s$-edge matching in the colored multigraph: every such matching has $s!2^s$ ordered orientations. Consequently this auxiliary law conditioned on $V$ is within $O(n^{-1})$ total variation of the original colored-union law with the uniform edge-selection kernel. Only the marginal on $G$ is reweighted, by $\Prob(V\mid G)=1+O(n^{-1})$.

Condition now on the sampled colors, all matchings with a touched color, and the sampled vertices. The event $V$ is determined by this data, so all untouched matchings remain independent uniform perfect matchings. Their number $r$ is at least $k-s=k/2\ge3$ for every color pattern. After surgery, the conditional internal adjacency matrix is
\begin{equation}\label{eq:border}
 \begin{pmatrix}M+F&u\\u^T&0\end{pmatrix},\qquad
 F\one=(k-r)\one-u,
\end{equation}
where $M$ is a sum of $r\ge3$ independent uniform matchings, $F$ is fixed and symmetric, and $u$ indicates exactly $k$ old vertices. The two internal blocks retain independent untouched matchings after their own exposed data are fixed.

\section{Conditional finite-field estimate}
Fix an odd prime $p\nmid\ell$. For every $r\in[k/2,k]$, use the established M\'esz\'aros mixing theorem \cite[Theorem 1.6]{meszaros}. Its reference mass for $Mx=t$ is
\[
 q_x^{(r)}(t)=\begin{cases}
 \boldsymbol1_{\{t=ra\one\}},&x=a\one,\\
 p^{-n+1}\boldsymbol1_{\{\one^Tt=r\one^Tx\}},&x\text{ nonconstant}.
 \end{cases}
\]
Let $\epsilon_n$ be the maximum, over these finitely many $r$, of the sum over $x$ of maximum-atom errors. Then $\epsilon_n\to0$ for every fixed $p$. The tensor restriction in that theorem is vacuous over odd $\Fp$, as in the companion proof.

Expose two internal blocks of form \eqref{eq:border}, denoting their stars by $u,z$, fixed matrices by $F_1,F_2$, and matching sums by $M_1,M_2$. Split the prescribed cross matrix at the two root vertices:
\[
 B=\begin{pmatrix}B_0&v\\w^T&b\end{pmatrix},\qquad b\in\{0,1\}.
\]
In particular $B_0\one=\ell\one-v$, $B_0^T\one=\ell\one-w$, and $\one^Tv=\one^Tw=\ell-b$.
Write an eigenvector as $(x,\alpha;y,\beta)$. The old-vertex equations are
\begin{align*}
 M_1x&=(\delta I-F_1)x-u\alpha-B_0y-v\beta,\\
 M_2y&=(\delta I-F_2)y-z\beta-B_0^Tx-w\alpha,
\end{align*}
and the two root equations are
\begin{align}
 u^Tx+w^Ty+b\beta&=\delta\alpha,\label{eq:r1}\\
 v^Tx+z^Ty+b\alpha&=\delta\beta.\label{eq:r2}
\end{align}

\paragraph{Affine bounds and total replacement error.}
For every fixed $c$, the reference sum
\[
 \sum_xq_x^{(r_1)}((\delta I-F_1)x+c)
\]
is at most two. For nonconstant $x$, its scalar constraint is
$(\ell\one-u)^Tx=\one^Tc$, a nonzero linear functional since $p\nmid\ell$ and the support of $u$ has bounded size. At most $p^{n-1}$ inputs each contribute $p^{-n+1}$. Constant $x=a\one$ requires $c=a(\ell\one-u)$, allowing at most one $a$. The actual affine sum is therefore at most $2+\epsilon_n$; the other block obeys the same bounds.

Expand the expected finite-field kernel size over $(x,y,\alpha,\beta)$, retaining the indicator of \eqref{eq:r1}--\eqref{eq:r2}. Conditional independence factors the old-vertex probabilities. For each fixed $(\alpha,\beta)$, replacing the first factor costs at most $(2+\epsilon_n)\epsilon_n$ and replacing the second costs at most $2\epsilon_n$, by these affine bounds. Dropping root indicators only enlarges this error bound. Total error is at most
\begin{equation}\label{eq:error}
 p^2(4\epsilon_n+\epsilon_n^2).
\end{equation}

\paragraph{Both old vectors nonconstant.}
After substitution of \eqref{eq:r1}, the first reference hyperplane constraint becomes
\begin{equation}\label{eq:r0}
 \one^Tx+\one^Ty+\alpha+\beta=0.
\end{equation}
For example, before substitution it is
\[
 \ell(\one^Tx+\one^Ty)-u^Tx-w^Ty+k\alpha+(\ell-b)\beta=0.
\]
The second hyperplane constraint similarly reduces to \eqref{eq:r0} after \eqref{eq:r2}. For each fixed root pair the three coefficient rows are
\[
 (\one^T,\one^T),\quad (u^T,w^T),\quad(v^T,z^T).
\]
They are independent for all sufficiently large $n$: a coordinate outside $\operatorname{supp}u\cup\operatorname{supp}v$ kills the first coefficient; one in $\operatorname{supp}u\setminus\operatorname{supp}v$ kills the second, since $k>\ell$; and $z\ne0$ kills the third. There are at most $p^{2n-3}$ old-vector pairs per root pair. Multiplying by $p^2$ root choices and $p^{-2n+2}$ reference masses bounds this contribution by $p$.

\paragraph{Both old vectors constant.}
If $x=a\one,y=c\one$, the first old vector equation is
\[
 \ell(a+c)\one+(\alpha-a)u+(\beta-c)v=0.
\]
Outside the local supports it forces $a+c=0$; on $\operatorname{supp}u\setminus\operatorname{supp}v$ it forces $\alpha=a$. The second old equation similarly forces $\beta=c$. These planted solutions satisfy the root equations, giving exactly $p$ choices.

\paragraph{Exactly one old vector constant.}
If $x=a\one$, the first old equation requires
\[
 B_0y=-\ell a\one-(\alpha-a)u-v\beta.
\]
For each of $p^3$ scalar choices $(a,\alpha,\beta)$ there are at most $p^{n-\rank B_0}$ choices of $y$, each with remaining mass at most $p^{-n+1}$. Including both orientations gives at most $2p^{4-\rank B_0}$. The deterministic disjoint-row-support argument gives
\[
 \rank_{\Fp}B_0\ge\rank_{\Fp}B-2
 \ge\frac{N}{\ell(\ell-1)+1}-2.
\]
Combining these cases with \eqref{eq:error}, uniformly over all exposed data and prescribed cross matrices,
\begin{equation}\label{eq:finalmoment}
 \E|\ker_{\Fp}(A_B-\delta I)|
 \le2p+2p^{4-\rank B_0}+p^2(4\epsilon_n+\epsilon_n^2).
\end{equation}
Average over the exposed data and take $n\to\infty$ for fixed $p$. Real nullity at least two forces finite-field nullity at least two, since all integer minors of order $2N-1$ vanish. Thus the limiting failure probability is at most $2/p$. Letting fixed admissible primes tend to infinity proves high-probability simplicity under the colored auxiliary surgery law.

\section{Order of the final model transfer}
The mixing calculation was applied before conditioning the original colored union on simplicity; that conditioning would destroy the independent untouched matchings. First transfer from the auxiliary experiment to the original colored-union law with uniform selected $E$ by their $O(n^{-1})$ total variation comparison. Next condition the original unions $G$ in the two blocks on simplicity, which has probability bounded below at fixed $k$. This costs only a bounded factor in the already-proved failure probability. On simple $G$, colored copies are ordinary edges, so the selection kernel is exactly the uncolored uniform $s$-edge matching kernel.

The established even-size contiguity between simple matching unions and uniform simple regular graphs \cite[Theorem 3]{mrrw} now transfers high-probability success with this same kernel. Contiguity survives an attached common Markov kernel: its conditional bad-event probability is a bounded function of $G$; Markov's inequality, contiguity, and boundedness transfer its vanishing expectation. Finite products handle the two independent blocks. The first section's $O(n^{-1})$ total variation estimate then yields independent uniform simple odd-size internal graphs. Fixing any deterministic sequence of cross matrices throughout and then choosing a violating maximizing sequence proves the uniformity in Theorem~\ref{thm:odd}.

For $k=4$, the two selected surgery edges can touch two colors, leaving only two independent untouched matchings. This occurs with nonvanishing probability. The imported degree-at-least-three mixing theorem therefore does not apply to that component; this proof does not establish the odd-size $k=4$ case. The qualitative mixing and contiguity steps also provide no effective convergence rate.

\begin{thebibliography}{9}
\bibitem{meszaros} A.~M\'esz\'aros, \emph{The distribution of sandpile groups of random regular graphs}, arXiv:1806.03736v4, 2020, Theorem~1.6.
\url{https://arxiv.org/pdf/1806.03736v4}.
\bibitem{mrrw} M.~S.~O.~Molloy, H.~Robalewska, R.~W.~Robinson, N.~C.~Wormald, \emph{1-Factorizations of random regular graphs}, Random Structures \& Algorithms \textbf{10} (1997), 305--321, Theorem~3. See also the explicit model chain in Nguyen--Wood, arXiv:1806.10068v2, Section~4, PDF p.~8.
\url{https://arxiv.org/pdf/1806.10068v2}.
\bibitem{hlmpw} L.~Hollom, L.~Lichev, A.~Mond, J.~Portier, Y.~Wang, \emph{Approximate Itai--Zehavi Conjecture for Random Graphs}, Random Structures \& Algorithms \textbf{69} (2026), e70095, odd-order regular-graph argument.
\url{https://onlinelibrary.wiley.com/doi/full/10.1002/rsa.70095}.
\end{thebibliography}
\end{document}
